\footnotesize
\begin{longtable}{>{\raggedright\arraybackslash}p{0.17\linewidth}>{\raggedright\arraybackslash}p{0.22\linewidth}>{\raggedright\arraybackslash}p{0.55\linewidth}}
\caption{Every modelling assumption in the CFD models (all other numbers reference sourced parameters).}\label{tab:assumptions}\\
\toprule Model & Parameter & Reason \\ \midrule \endhead
aph & obstacles[0]\allowbreak{}.\allowbreak{}center & Science Carrier fills the chamber floor ('installed into the base of the growth chamber', registry 2017); footprint equal to the chamber footprint by construction (enclosure footprint was derived from the carrier). \\
aph & zones[0]\allowbreak{}.\allowbreak{}size & Canopy zone spans the planted surface minus 10 mm margins; 100 mm height represents a young canopy (PH-04 plants were 14 days old at the first RH change). Sensitivity: canopy height not sourced. \\
aph & zones[0]\allowbreak{}.\allowbreak{}center & Canopy zone sits directly on the Science Carrier surface (z = 51 mm). \\
aph & openings[0]\allowbreak{}.\allowbreak{}size & Vent sizes not published ('Vent sizes not published', aph.yaml). 300 x 60 mm slots assumed on each ECS side. \\
aph & openings[0]\allowbreak{}.\allowbreak{}center & ECS returns air 'at plant level from opposing sides' (Monje et al. 2020): inlets centred 50 mm above the carrier. \\
aph & openings[1]\allowbreak{}.\allowbreak{}size & As openings[0]. \\
aph & openings[1]\allowbreak{}.\allowbreak{}center & As openings[0]. \\
aph & openings[2]\allowbreak{}.\allowbreak{}size & Outlet size not published; equal to inlets. \\
aph & openings[2]\allowbreak{}.\allowbreak{}center & ECS fans 'remove air from the top of the growth chamber above the plants' (Monje et al. 2020): outlets centred 40 mm below the ceiling on the same side walls (ECS modules are 'mounted on each side'). \\
aph & openings[3]\allowbreak{}.\allowbreak{}size & As openings[2]. \\
aph & openings[3]\allowbreak{}.\allowbreak{}center & As openings[2]. \\
aph & openings.\allowbreak{}inlet\_\allowbreak{}velocity\_\allowbreak{}interpretation & The PH-04 'wind speed of 0.3 m/s' set-point is applied as the inlet face velocity; where APH measures wind speed is not stated (aph.yaml note). \\
aph & zones.\allowbreak{}canopy.\allowbreak{}co2\_\allowbreak{}load & Canopy net CO2 uptake uses the wheat range 7-25 umol m-2 s-1 (Wheeler 1996 as cited by Monje et al. 2020; mid-point 16) over the 0.19 m2 growth area; no pepper rate was found. Transpiration and lamp/canopy heat are not modelled (no sourced value). \\
aph\_\allowbreak{}fanmax & obstacles[0]\allowbreak{}.\allowbreak{}center & Science Carrier fills the chamber floor ('installed into the base of the growth chamber', registry 2017); footprint equal to the chamber footprint by construction (enclosure footprint was derived from the carrier). \\
aph\_\allowbreak{}fanmax & zones[0]\allowbreak{}.\allowbreak{}size & Canopy zone spans the planted surface minus 10 mm margins; 100 mm height represents a young canopy (PH-04 plants were 14 days old at the first RH change). Sensitivity: canopy height not sourced. \\
aph\_\allowbreak{}fanmax & zones[0]\allowbreak{}.\allowbreak{}center & Canopy zone sits directly on the Science Carrier surface (z = 51 mm). \\
aph\_\allowbreak{}fanmax & openings[0]\allowbreak{}.\allowbreak{}size & Vent sizes not published ('Vent sizes not published', aph.yaml). 300 x 60 mm slots assumed on each ECS side. \\
aph\_\allowbreak{}fanmax & openings[0]\allowbreak{}.\allowbreak{}center & ECS returns air 'at plant level from opposing sides' (Monje et al. 2020): inlets centred 50 mm above the carrier. \\
aph\_\allowbreak{}fanmax & openings[1]\allowbreak{}.\allowbreak{}size & As openings[0]. \\
aph\_\allowbreak{}fanmax & openings[1]\allowbreak{}.\allowbreak{}center & As openings[0]. \\
aph\_\allowbreak{}fanmax & openings[2]\allowbreak{}.\allowbreak{}size & Outlet size not published; equal to inlets. \\
aph\_\allowbreak{}fanmax & openings[2]\allowbreak{}.\allowbreak{}center & ECS fans 'remove air from the top of the growth chamber above the plants' (Monje et al. 2020): outlets centred 40 mm below the ceiling on the same side walls (ECS modules are 'mounted on each side'). \\
aph\_\allowbreak{}fanmax & openings[3]\allowbreak{}.\allowbreak{}size & As openings[2]. \\
aph\_\allowbreak{}fanmax & openings[3]\allowbreak{}.\allowbreak{}center & As openings[2]. \\
aph\_\allowbreak{}fanmax & openings.\allowbreak{}inlet\_\allowbreak{}velocity\_\allowbreak{}interpretation & The PH-04 'wind speed of 0.3 m/s' set-point is applied as the inlet face velocity; where APH measures wind speed is not stated (aph.yaml note). \\
aph\_\allowbreak{}fanmax & zones.\allowbreak{}canopy.\allowbreak{}co2\_\allowbreak{}load & Canopy net CO2 uptake uses the wheat range 7-25 umol m-2 s-1 (Wheeler 1996 as cited by Monje et al. 2020; mid-point 16) over the 0.19 m2 growth area; no pepper rate was found. Transpiration and lamp/canopy heat are not modelled (no sourced value). \\
aqh & openings[0]\allowbreak{}.\allowbreak{}size & Water inlet/outlet geometry of the aquarium is not published (aqh.yaml: 'aquarium water inlet/outlet geometry'). 20 x 20 mm ports assumed. \\
aqh & openings[0]\allowbreak{}.\allowbreak{}center & Inlet low on one end wall, outlet high on the opposite end wall (assumed cross-flow arrangement). \\
aqh & openings[1]\allowbreak{}.\allowbreak{}size & As openings[0]. \\
aqh & openings[1]\allowbreak{}.\allowbreak{}center & As openings[0]. \\
aqh & domain.\allowbreak{}water\_\allowbreak{}only & The aquarium is modelled full of water; the air pocket of the air-water interface ('air stabilizer') and the fish are not represented (no dimensions or O2 uptake sourced). Walls held at the set-point temperature; no heat sources, so 0 g and 1 g are expected to coincide. \\
mds & zones[0]\allowbreak{}.\allowbreak{}size & Mouse represented as a distributed source volume of 80 x 30 x 30 mm (adult-mouse scale) resting on the cage floor; body geometry not modelled as a solid. \\
mds & zones[0]\allowbreak{}.\allowbreak{}center & Mouse at the floor centre. \\
mds & openings[0]\allowbreak{}.\allowbreak{}size & Flow direction through the cage is not stated; waste filters are 'located below each Habitat' (Liu 2009, mds.yaml note), so air is taken to enter through the whole top grid and leave through the whole floor grid. \\
mds & openings[0]\allowbreak{}.\allowbreak{}center & As openings[0].size. \\
mds & openings[1]\allowbreak{}.\allowbreak{}size & As openings[0].size. \\
mds & openings[1]\allowbreak{}.\allowbreak{}center & As openings[0].size. \\
mds & openings.\allowbreak{}velocity\_\allowbreak{}interpretation & Cancedda 2012: airflow of 0.2-0.3 m/s through the cages; applied as a uniform face velocity (mid-point 0.25 m/s). \\
mds & environment.\allowbreak{}co2\_\allowbreak{}reference & In-cage CO2 is controlled but not reported; supply air is set to the ISS U.S. Lab cabin CO2 measured during RR-1 (Choi et al. 2020; mid-point of 788-4690 ppm). MDS flew in 2009 in JEM, so this is a proxy. \\
mds & zones.\allowbreak{}mouse.\allowbreak{}co2\_\allowbreak{}load & Organism-level C57BL/6 rates (ICES-2014-101) and BALB/c heat and moisture (Riskowski 2006) scaled by the generic 30 g adult mass from Cancedda 2012; MDS flew C57BL/10J males. \\
vented\_\allowbreak{}fly\_\allowbreak{}box\_\allowbreak{}vent20 & obstacles[0]\allowbreak{}.\allowbreak{}center & 15 vials (24 x 75 mm) fit the 94 x 94 x 130 mm box only lying along x in a 3 x 5 stack (upright, at most 12 fit); g-load at launch was directed into the food at the vial bottoms (OSD-347). Vials treated as solid (food + flies inside not resolved). \\
vented\_\allowbreak{}fly\_\allowbreak{}box\_\allowbreak{}vent20 & obstacles[0]\allowbreak{}.\allowbreak{}repeat.\allowbreak{}n & 3 x 5 stack, see centre. \\
vented\_\allowbreak{}fly\_\allowbreak{}box\_\allowbreak{}vent20 & obstacles[0]\allowbreak{}.\allowbreak{}repeat.\allowbreak{}pitch & Vials spaced 32 mm (y) and 25 mm (z) to fill the box cross-section. \\
vented\_\allowbreak{}fly\_\allowbreak{}box\_\allowbreak{}vent20 & zones[0]\allowbreak{}.\allowbreak{}H & CO2 leaves each vial through its open foam plug (OSD-347; Researcher's Guide); released in a 4 mm layer at the plug end (x = 84.5-88.5 mm). \\
vented\_\allowbreak{}fly\_\allowbreak{}box\_\allowbreak{}vent20 & zones[0]\allowbreak{}.\allowbreak{}center & Plug end of each vial. \\
vented\_\allowbreak{}fly\_\allowbreak{}box\_\allowbreak{}vent20 & zones[0]\allowbreak{}.\allowbreak{}repeat.\allowbreak{}n & As obstacles. \\
vented\_\allowbreak{}fly\_\allowbreak{}box\_\allowbreak{}vent20 & zones[0]\allowbreak{}.\allowbreak{}repeat.\allowbreak{}pitch & As obstacles. \\
vented\_\allowbreak{}fly\_\allowbreak{}box\_\allowbreak{}vent20 & zones.\allowbreak{}plug.\allowbreak{}sources.\allowbreak{}CO2.\allowbreak{}count & 13 flies per vial (OSD-347: '8 virgin females and 5 male adult virgin flies') x 15 vials = 195 adults; progeny that emerge during the mission (20-60 per vial on return) are not included, so the load is a lower bound. \\
vented\_\allowbreak{}fly\_\allowbreak{}box\_\allowbreak{}vent20 & openings[0]\allowbreak{}.\allowbreak{}size & Vent hole number, size and position are not published (vented\_fly\_box.yaml). Run at two sizes, 20 x 20 and 60 x 60 mm per end wall, to show the sensitivity. \\
vented\_\allowbreak{}fly\_\allowbreak{}box\_\allowbreak{}vent20 & openings[0]\allowbreak{}.\allowbreak{}center & Vents centred on the end walls. \\
vented\_\allowbreak{}fly\_\allowbreak{}box\_\allowbreak{}vent20 & openings[1]\allowbreak{}.\allowbreak{}size & As openings[0]. \\
vented\_\allowbreak{}fly\_\allowbreak{}box\_\allowbreak{}vent20 & openings[1]\allowbreak{}.\allowbreak{}center & As openings[0]. \\
vented\_\allowbreak{}fly\_\allowbreak{}box\_\allowbreak{}vent20 & openings.\allowbreak{}exterior & The exterior at each vent is taken as well mixed ISS cabin air (CO2 mid-point of the RR-1 U.S. Lab range, Choi et al. 2020) at the box temperature. \\
vented\_\allowbreak{}fly\_\allowbreak{}box\_\allowbreak{}vent60 & obstacles[0]\allowbreak{}.\allowbreak{}center & 15 vials (24 x 75 mm) fit the 94 x 94 x 130 mm box only lying along x in a 3 x 5 stack (upright, at most 12 fit); g-load at launch was directed into the food at the vial bottoms (OSD-347). Vials treated as solid (food + flies inside not resolved). \\
vented\_\allowbreak{}fly\_\allowbreak{}box\_\allowbreak{}vent60 & obstacles[0]\allowbreak{}.\allowbreak{}repeat.\allowbreak{}n & 3 x 5 stack, see centre. \\
vented\_\allowbreak{}fly\_\allowbreak{}box\_\allowbreak{}vent60 & obstacles[0]\allowbreak{}.\allowbreak{}repeat.\allowbreak{}pitch & Vials spaced 32 mm (y) and 25 mm (z) to fill the box cross-section. \\
vented\_\allowbreak{}fly\_\allowbreak{}box\_\allowbreak{}vent60 & zones[0]\allowbreak{}.\allowbreak{}H & CO2 leaves each vial through its open foam plug (OSD-347; Researcher's Guide); released in a 4 mm layer at the plug end (x = 84.5-88.5 mm). \\
vented\_\allowbreak{}fly\_\allowbreak{}box\_\allowbreak{}vent60 & zones[0]\allowbreak{}.\allowbreak{}center & Plug end of each vial. \\
vented\_\allowbreak{}fly\_\allowbreak{}box\_\allowbreak{}vent60 & zones[0]\allowbreak{}.\allowbreak{}repeat.\allowbreak{}n & As obstacles. \\
vented\_\allowbreak{}fly\_\allowbreak{}box\_\allowbreak{}vent60 & zones[0]\allowbreak{}.\allowbreak{}repeat.\allowbreak{}pitch & As obstacles. \\
vented\_\allowbreak{}fly\_\allowbreak{}box\_\allowbreak{}vent60 & zones.\allowbreak{}plug.\allowbreak{}sources.\allowbreak{}CO2.\allowbreak{}count & 13 flies per vial (OSD-347: '8 virgin females and 5 male adult virgin flies') x 15 vials = 195 adults; progeny that emerge during the mission (20-60 per vial on return) are not included, so the load is a lower bound. \\
vented\_\allowbreak{}fly\_\allowbreak{}box\_\allowbreak{}vent60 & openings[0]\allowbreak{}.\allowbreak{}size & Vent hole number, size and position are not published (vented\_fly\_box.yaml). Run at two sizes, 20 x 20 and 60 x 60 mm per end wall, to show the sensitivity. \\
vented\_\allowbreak{}fly\_\allowbreak{}box\_\allowbreak{}vent60 & openings[0]\allowbreak{}.\allowbreak{}center & Vents centred on the end walls. \\
vented\_\allowbreak{}fly\_\allowbreak{}box\_\allowbreak{}vent60 & openings[1]\allowbreak{}.\allowbreak{}size & As openings[0]. \\
vented\_\allowbreak{}fly\_\allowbreak{}box\_\allowbreak{}vent60 & openings[1]\allowbreak{}.\allowbreak{}center & As openings[0]. \\
vented\_\allowbreak{}fly\_\allowbreak{}box\_\allowbreak{}vent60 & openings.\allowbreak{}exterior & The exterior at each vent is taken as well mixed ISS cabin air (CO2 mid-point of the RR-1 U.S. Lab range, Choi et al. 2020) at the box temperature. \\
\bottomrule
\end{longtable}
\normalsize
